% from ega4-2-fr.tex
\oldpage[IV-2]{108}

nilradical de $X$ étant un $\mathcal O_X$-Module cohérent (I, 6.1.1 et
$0_{\mathrm I}$, 5.2.2)), il revient au même de dire que $X$ possède la
propriété $(R_0)$ ou que l'ensemble $U$ est partout dense. Par suite:
\end{env}

\begin{proposition}[5.8.5]
\label{IV.2.s5.8.5-fr}
Pour qu'un préschéma localement noethérien $X$ soit réduit, il faut et il
suffit qu'il vérifie les propriétés $(S_1)$ et $(R_0)$.
\end{proposition}

Compte tenu de (5.7.5), cela résulte de la remarque précédente et de (3.2.1).

\begin{theorem}[5.8.6]
\EGAReaderTarget{ega4.c04.thm.5-8-6}{IV.2.s5.8.6-fr}
(critère de Serre). Soit $X$ un préschéma localement noethérien. Pour que
$X$ soit normal, il faut et il suffit que $X$ vérifie les propriétés $(S_2)$
et $(R_1)$, autrement dit, que pour tout $x\in X$, on ait les propriétés
suivantes:
\begin{enumerate}
\item[(i)] Si $\dim(\mathcal O_x)\leq1$, $\mathcal O_x$ est régulier
(c'est-à-dire est un corps ou un anneau de valuation discrète (0, 17.1.4)).
\item[(ii)] Si $\dim(\mathcal O_x)\geq2$, alors
$\operatorname{prof}(\mathcal O_x)\geq2$.
\end{enumerate}
\end{theorem}

Les conditions sont nécessaires. En effet, dire que $X$ est normal signifie
que pour tout $x\in X$, $\mathcal O_x$ est un anneau local noethérien
intégralement clos. Si $\dim(\mathcal O_x)=0$ (resp.
$\dim(\mathcal O_x)=1$), on en conclut que $\mathcal O_x$ est un corps
puisque $\mathcal O_x$ est intègre (resp. que $\mathcal O_x$ est un anneau
de valuation discrète, en vertu de (II, 7.1.6)). D'autre part, pour tout
élément $f_x\neq0$ de $\mathcal O_x$, on sait (Bourbaki, \emph{Alg. comm.},
chap. VII, § 1, n° 4, prop. 8) que les idéaux premiers associés à
$\mathcal O_x/f_x\mathcal O_x$ sont non immergés, donc $\mathcal O_x$
vérifie $(S_2)$ (5.7.7).

Les conditions sont suffisantes. En effet, il résulte d'abord de (5.8.5) que
$X$ est réduit. La question étant locale, on peut en outre supposer que
$X=\operatorname{Spec}(A)$, où $A$ est un anneau noethérien réduit
(I, 5.1.4); si $R$ est l'anneau total des fractions de $A$, $R$ est composé
direct d'un nombre fini de corps, et (compte tenu de (II, 6.3.6)), il suffira
de prouver que $A$ est intégralement fermé dans $R$. Soit donc $h=f/g$ un
élément de $R$ entier sur $A$, $g$ et $f$ étant des éléments de $A$ tels que
$g$ soit non diviseur de 0. On a une relation de la forme
\begin{equation}
f^n+\sum_{i=1}^n a_i f^{n-i}g^i=0
\qquad\text{avec }a_i\in A\quad\text{pour }1\leq i\leq n.
\tag{5.8.6.1}\label{IV.2.s5.8.6.1-fr}
\end{equation}

Soit $\mathfrak p$ un idéal premier de $A$ tel que
$\dim(A_{\mathfrak p})=1$; si $f_{\mathfrak p}$ et $g_{\mathfrak p}$ sont
les images de $f$ et $g$ dans $A_{\mathfrak p}$, il résulte de (5.8.6.1) que
$f_{\mathfrak p}/g_{\mathfrak p}$ (qui appartient à l'anneau total des
fractions de $A_{\mathfrak p}$, puisque $g_{\mathfrak p}$ est non diviseur
de 0 dans $A_{\mathfrak p}$ par platitude $(0_{\mathrm I},5.3.1)$) est entier
sur $A_{\mathfrak p}$; mais comme $A_{\mathfrak p}$ est régulier, donc
intégralement clos, on a $f_{\mathfrak p}/g_{\mathfrak p}\in A_{\mathfrak p}$.
En d'autres termes, on a
$(fA)_{\mathfrak p}\subset(gA)_{\mathfrak p}$. Mais l'hypothèse $(S_2)$
entraîne (5.7.7), puisque $g$ est non diviseur de 0 dans $A$, que $A/gA$
n'a que des idéaux premiers associés non immergés $\mathfrak p_i$
($1\leq i\leq n$); or, $gA$ est l'intersection d'idéaux primaires
$\mathfrak q_i$ correspondant aux $\mathfrak p_i$, et d'après ce qu'on vient
de voir, les $\mathfrak q_i$ sont les images réciproques dans $A$, par les
homomorphismes $A\to A_{\mathfrak p_i}$, des idéaux
$(gA)_{\mathfrak p_i}$ (Bourbaki, \emph{Alg. comm.}, chap. IV, § 2, n° 3,
prop. 5). Mais en vertu du Hauptidealsatz (0, 16.3.2) on a
$\dim(A_{\mathfrak p_i})=1$ pour $1\leq i\leq n$, donc
$(fA)_{\mathfrak p_i}\subset(gA)_{\mathfrak p_i}$ pour tout $i$ d'après ce
qui précède; comme $fA$ est contenu dans l'intersection des images
réciproques des $(fA)_{\mathfrak p_i}$ ($1\leq i\leq n$), on a
$fA\subset gA$, c'est-à-dire $f/g\in A$. C.Q.F.D.

